\subsection{Characterization of completely secant ideals}

Let A be a ring and J an ideal of A. The graded A-module \(\bigoplus_{n\in\mathbf{N}}J^{n}\) has a natural structure of graded A-algebra, deduced from the multiplication in the ring A; the identity map of J into \(J^{1}\) therefore extends to a surjective homomorphism of graded A-algebras, said to be canonical

\[
\alpha_ {J}: \mathsf {S} _ {A} (J) \longrightarrow \bigoplus_ {n \in N} J ^ {n}.
\]

By extension of scalars to the ring A/J, from \(\alpha_{\mathrm{J}}\) one deduces a surjective homomorphism of graded A/J-algebras, also said to be canonical

\[
\beta_ {J}: \mathsf {S} _ {\mathrm{A} / J} (\mathrm{J} / \mathrm{J} ^ {2}) \longrightarrow \operatorname{gr} _ {J} (\mathrm{A}),
\]

\(\operatorname{avec} \mathrm{gr}_{\mathrm{J}}(\mathrm{A}) = \bigoplus_{n \in \mathbf{N}} \mathrm{J}^n / \mathrm{J}^{n+1}\) .

THEOREM 1.--- Let A be a Noetherian ring and J an ideal of A. The following conditions are equivalent:

\begin{enumerate}
\def\labelenumi{(\roman{enumi})}
\item
  the ideal J is completely secant ;
\item
  the ideal J is completely secant at every maximal ideal \(\mathfrak{m} \in V(J)\) ;
\item
  the A/J-module J/J² is projective and the canonical homomorphism \(\alpha_{J}: S_{\Lambda}(J) \longrightarrow \bigoplus_{n \in \mathbb{N}} J^{n}\) is bijective;
\item
  the \(\Lambda/\mathrm{J}\) -module \(\mathrm{J}/\mathrm{J}^2\) is projective and the canonical homomorphism \(\beta_{\mathrm{J}}:\mathbf{S}_{\mathrm{A}/\mathrm{J}}(\mathrm{J}/\mathrm{J}^2)\longrightarrow\mathrm{gr}_{\mathrm{J}}(\mathrm{A})\) is bijective.
\end{enumerate}

(i)⇒(ii): this is trivial.

(ii)⇒(iii): suppose condition (ii) is satisfied. It suffices to prove that for every maximal ideal m of A, the \(A_{m}/J_{m}\) -module \(J_{m}/J_{m}^{2}\) is free and the homomorphism \(\alpha_{J_{m}}: S_{A_{m}}(J_{m}) \longrightarrow \bigoplus_{n} J_{m}^{n}\) is bijective (II, § 5, no. 2, th. 1 and § 3, no. 3, th. 1). But these assertions are trivial when m does not belong to V(J), since then we have \(J_{m} = A_{m}\) and they follow from A, X, p.~160, th. 1 and p.~161, remark, when m belongs to V(J).

\begin{enumerate}
\def\labelenumi{(\roman{enumi})}
\setcounter{enumi}{2}
\tightlist
\item
  \(\Rightarrow\) (iv): this is clear.
\end{enumerate}

\((\mathrm{iv}) \Rightarrow (\mathrm{i})\): Suppose condition (iv) is satisfied; let \(\mathfrak{p}\) be a prime ideal of A containing J. Then the \(A_{\mathfrak{p}} / J_{\mathfrak{p}}\) -module \(J_{\mathfrak{p}} / J_{\mathfrak{p}}^2\) is free. Let \(\mathbf{x} = (x_1, \ldots, x_r)\) be a sequence of elements of \(J_{\mathfrak{p}}\) lifting a basis of \(J_{\mathfrak{p}} / J_{\mathfrak{p}}^2\) . The sequence \(\mathbf{x}\) generate \(J_{\mathfrak{p}}\) (Nakayama's lemma), and satisfies by construction condition (iv) of prop. 1. Consequently the ideal \(J_{\mathfrak{p}}\) of \(A_{\mathfrak{p}}\) is completely secant, and J is completely secant at \(\mathfrak{p}\) .

Remark 1. Suppose the ideal J is completely secant; let \((x_1, \ldots, x_r)\) be a sequence of elements of J, such that for every maximal ideal \(\mathfrak{m} \in V(J)\) the canonical images of the \(x_i\) in J/mJ form a basis of this A/m-vector space. Then the A/J-module \(J/J^2\) is free and the canonical images in \(J/J^2\) of the \(x_i\) form a basis of it: it suffices indeed to verify that the images of the \(x_i\) form a basis of the \(A_{\mathfrak{m}}/J_{\mathfrak{m}}\) -module \(J_{\mathfrak{m}}/J_{\mathfrak{m}}^2\) for every \(\mathfrak{m} \in V(J)\) (II, § 3, n° 3, th. 1), which follows from loc. cit., n° 2, prop. 5 and cor. 2, since the \(A_{\mathfrak{m}}/J_{\mathfrak{m}}\) -module \(J_{\mathfrak{m}}/J_{\mathfrak{m}}^2\) is projective (th. 1).

COROLLARY 1.--- Let \(\widehat{\mathrm{A}}\) the separated completion of A for the J-adic topology and \(\widehat{\mathrm{J}} = \mathrm{J}\widehat{\mathrm{A}}\) the separated completion of J. For the ideal \(\widehat{\mathrm{J}}\) of \(\widehat{\mathrm{A}}\) to be completely secant, it is necessary and sufficient that the ideal J of A be completely secant.

Indeed, the canonical map \(\mathrm{gr}_{\mathrm{J}}(\mathsf{A}) \to \mathrm{gr}_{\widehat{\mathrm{J}}}(\widehat{\mathsf{A}})\) is an isomorphism of graded rings; it therefore suffices to apply criterion (iv).

More generally:

COROLLARY 2.--- Let ρ : A → B be a homomorphism of Noetherian rings making B a flat A-module, and let J be an ideal of A.

\begin{enumerate}
\def\labelenumi{\alph{enumi})}
\item
  If J is completely secant, the ideal JB of B is completely secant.
\item
  Suppose that the ideal JB of B is completely secant and that every maximal ideal \(\mathfrak{m} \in \mathrm{V}(\mathrm{J})\) is the inverse image of a maximal ideal of B. Then the ideal J is completely secant. This is the case for example if B is a faithfully flat A-module.
\end{enumerate}

First, note that since the B is a flat A-module, \(J^n \otimes_A B\) is identified with \(J^n B\) and \((J^n / J^{n+1}) \otimes_{A/J} (B/JB)\) is identified with \(J^n B / J^{n+1} B\) for every integer \(n \geqslant 0\) . Assertion a) then follows from criterion (iii). Under the hypotheses of b), the A/J-module B/JB is faithfully flat (I, § 2, no. 7, cor. 2 of prop. 8 and § 3, no. 5, prop. 9) and J is completely secant according to criterion (iv), taking into account I, § 3, no. 1, prop. 2 and no. 6, prop. 12. The last assertion follows from prop. 8 of I, § 3, no. 5.

COROLLARY 3.-- Let A be a Noetherian ring and J a completely secant ideal of A. If A is a Macaulay ring (resp. a Gorenstein ring), then so is A/J.

Indeed, let m be a maximal ideal of A containing J. The ideal \(J_{m}\) of \(\Lambda_{m}\) is generated by a sequence \(A_{m}\) -regular; therefore \((A/J)_{m}\) is a Macaulay ring (resp. a Gorenstein ring) by Example 6 of § 2, No. 5 (resp. Example 2 of § 3, No. 9).

Remark 2. — A Noetherian ring A is said to be a complete intersection ring if, for every maximal ideal m of A, the complete local ring \(\widehat{A_{m}}\) is isomorphic to the quotient of a regular complete local Noetherian ring by a complete intersection ideal. It follows from cor. 3 above and cor. 2 of prop. 12 of § 3, no. 8 that such a ring is a Gorenstein ring.

